\section{논의}

\subsection{이 논문이 설득력 있게 말할 수 있는 것}

현재 결과를 기준으로 가장 설득력 있는 논문은 ``새로운 low-bit SOTA'' 논문이 아니다. 대신 \emph{회전 기반 KV-cache 압축을 Lie 군 관점에서 재해석하고, 공분산 정렬 정적 회전으로서 PCA의 의미를 수학과 실험으로 정당화하는 논문}이다. 이 주장은 이미 확보된 결과와 잘 맞는다. 비등방성, facet separability, PCA의 강한 surrogate 성능, query dynamicity의 분해는 모두 같은 방향을 가리킨다.

\subsection{왜 GPT-2가 여전히 필요했는가}

GPT-2를 사용한 이유도 이 서사에서 분명해진다. GPT-2는 최신 SOTA 모델이 아니라 \emph{회전 선택 문제를 해부하기 좋은 해석용 모델}이다. head별 공분산, random/PCA/identity 비교, 연속 회전 최적화, query-dependent schedule 같은 실험은 GPT-2에서 빠르고 깨끗하게 수행된다. 즉 GPT-2는 최종 실용성의 대리물이 아니라, 이론의 메커니즘을 읽기 위한 실험실 역할을 한다.

\subsection{무엇이 아직 비어 있는가}

비어 있는 부분도 분명하다. 첫째, 현대 GQA 모델에서 PCA 기반 회전이 strong baseline을 실제로 넘어서는지 아직 불확실하다. 둘째, 현재 FOKVQ의 성능 병목이 회전 선택이 아니라 회전 후 양자화기 설계라는 해석은 그럴듯하지만, 이를 분리 검증하는 직접 실험이 더 필요하다. 셋째, Lie 군 해석을 복소 unitary, banded generator, commutator regularization 등으로 확장한 변형들은 아직 탐색 단계다.

\subsection{따라서 어떤 논문이 되어야 하는가}

이 초안이 노려야 할 최종 형태는 다음과 같다. 본문에서는 Lie 군 관점의 정당화와 현재까지의 실증을 메인으로 두고, low-bit 메인 비교표는 ``긍정적 증명''이 아니라 ``어디까지는 되고 어디서부터 막히는가''를 보여 주는 honesty section으로 둔다. 만약 이후 추가 실험에서 stronger quantizer가 얹혀 실제 우위가 확인되면, 그때 논문은 ``Lie-guided quantizer design''으로 확장될 수 있다. 현재 시점에서는 그 이전 단계가 가장 정직하다.
